\documentclass[12pt]{article}

\usepackage[english]{babel}
\usepackage{amsmath,amssymb,amsthm}
\usepackage{fullpage}
\usepackage[hidelinks]{hyperref}

\allowdisplaybreaks

\newtheoremstyle{plainsl}
  {\topsep}
  {\topsep}
  {\slshape}
  {}
  {\normalfont\bfseries}
  {.}
  { }
  {}

\theoremstyle{plainsl}
\newtheorem{theorem}{Theorem}[section]
\newtheorem{lemma}[theorem]{Lemma}
\newtheorem*{conjecture}{Conjecture}

\numberwithin{equation}{section}

\newcommand{\NOne}{N_1}
\newcommand{\NTwo}{N_2}

\title{A dense-case theorem for Seymour's second-neighborhood conjecture}
\author{Jake Brukhman}
\date{9 August 2026}

\hypersetup{
  pdftitle={A dense-case theorem for Seymour's second-neighborhood conjecture},
  pdfauthor={Jake Brukhman}
}

\begin{document}
\maketitle

\begingroup
\renewcommand\thefootnote{}
\footnotetext[0]{Email: \texttt{jake@coinfund.io}}
\endgroup

\begin{abstract}
Seymour's Second Neighborhood Conjecture asserts that every finite oriented
graph has a vertex with at least as many exact second outneighbors as
outneighbors.  Established cases include tournaments, proved by
Fisher~\cite{Fisher1996}, and oriented graphs of minimum outdegree at most six,
proved by Kaneko and Locke~\cite{KanekoLocke2001}; a recent preprint treats
minimum outdegree seven~\cite{SadhukhanSandeepSen2026}.  For dense incomplete
graphs, Fidler and Yuster proved the conjecture when the missing edges form a
matching, a star, or a clique~\cite{FidlerYuster2007}, and Ghazal extended this
direction to generalized stars~\cite{Ghazal2012}.  We give a short fixed-target
counting proof that the conjecture holds for every oriented graph of order $n$
and minimum outdegree $\delta$ satisfying $n\le 2\delta+2$, a dense-case with no prescribed
structure on the missing edges.

\end{abstract}

\section{Introduction}

A finite \emph{directed graph} is a pair
$G=(V(G),A(G))$, where $V(G)$ is a finite set of vertices and
$A(G)$ is a set of ordered pairs of distinct vertices, called \emph{arcs}.
We write
\[
                              u\to v
\]
when $(u,v)\in A(G)$.  Thus the arrow points from the source of the arc to
its target.  The directed graph is \emph{oriented} if, for every two distinct
vertices $u,v$, at most one of $u\to v$ and $v\to u$ is present. For a vertex $v\in V(G)$, define its outdegree and indegree by
\[
 d^+(v)=|\{x\in V(G):v\to x\}|,
 \qquad
 d^-(v)=|\{x\in V(G):x\to v\}|.
\]
The minimum outdegree of $G$ is
$\delta=\min_{v\in V(G)}d^+(v).$ For a vertex $v\in V(G)$, define the exact first outneighborhood by
\[
                    \NOne(v)=\{x\in V(G):v\to x\},
\]
and the exact second outneighborhood by
\[
 \NTwo(v)=
 \{x\in V(G)\setminus\NOne(v):
   y\to x\text{ for some }y\in\NOne(v)\}.
\]
Direct outneighbors are excluded explicitly.  The vertex $v$ is excluded
automatically: a path $v\to y\to v$ would be a directed $2$-cycle, which an
oriented graph cannot contain.  Thus $\NTwo(v)$ consists exactly of the
vertices reached from $v$ in two steps but not in one step.  In a reachability statement from $v$ to $x$, we
call $v$ the \emph{source vertex} and $x$ the \emph{target vertex}. A vertex $v$ is a \emph{Seymour vertex} if
$|\NTwo(v)|\ge|\NOne(v)|.$

\begin{conjecture}[Seymour]
Every finite oriented graph has a Seymour vertex.
\end{conjecture}

Several results bear directly on dense or extremal cases.  Fidler and
Yuster~\cite{FidlerYuster2007} proved the conjecture for several prescribed
missing-edge structures, and Ghazal~\cite{Ghazal2012} generalized that line of
work.  Espuny D\'iaz, Gir\~ao, Granet, and
Kronenberg~\cite{EspunyDiazEtAl2024} obtained reductions involving minimum
outdegree, and Guo, Kang, and Zwaneveld~\cite{GuoKangZwaneveld2026} studied
Seymour-tight orientations and counterexamples close to regular tournaments. We present the following result.

\begin{theorem}[Fixed-target theorem]\label{thm:fixed-target}
Let $G$ be an oriented graph on $n$ vertices with minimum outdegree $\delta$.
If
\[
                              n\le2\delta+2,
\]
then $G$ has a Seymour vertex.
\end{theorem}

\section{The fixed-target capacity bound}

For a source vertex $v$, define
\[
 U(v)=V(G)\setminus\bigl(\{v\}\cup\NOne(v)\cup\NTwo(v)\bigr).
\]
Thus $U(v)$ is the set of target vertices that the source $v$ cannot reach in
at most two steps.

Now fix a target vertex $x$.  Define the \emph{trap} of $x$ by
\[
 T(x)=\{y\in V(G)\setminus\{x\}:y\nrightarrow x\}.
\]
These are exactly the vertices that do not send an arc to $x$.  Because $G$
is oriented, every outneighbor of $x$ belongs to $T(x)$.  Hence
\[
                         |T(x)|\ge d^+(x)\ge\delta.
\]
Define the \emph{target slack}
\[
                         q(x)=|T(x)|-\delta\ge0.
\]
Here both $|T(x)|$ and $\delta$ count vertices: $\delta$ is the minimum
number of outneighbors at any vertex.  Thus $q(x)$ is the number of vertices
by which the trap exceeds the size of a minimum outneighborhood.

Define
\[
                         P(x)=\{v\in V(G):x\in U(v)\}.
\]
Thus $P(x)$ is the set of source vertices from which the fixed target $x$
cannot be reached in at most two steps.  If $v\in P(x)$, then $v\nrightarrow
x$.  Moreover, no vertex of $\NOne(v)$ points to $x$, since otherwise there
would be a two-step path from $v$ to $x$.  Therefore
\[
                         \{v\}\cup\NOne(v)\subseteq T(x).   \tag{1}
\]
This is the fixed-target observation: every source $v$ that cannot reach $x$ in 2 steps, together with
all of its outneighbors, is confined to the same trap $T(x)$.

\begin{lemma}[Fixed-target capacity]\label{lem:capacity}
For every target vertex $x$,
\[
 |P(x)|=0\quad\text{if }q(x)=0,
 \qquad
 |P(x)|\le2q(x)-1\quad\text{if }P(x)\ne\varnothing.
\]
\end{lemma}

\begin{proof}
Fix $x$ and set $p:=|P(x)|$.
By (1), $P(x)\subseteq T(x)$, and every arc whose source lies in $P(x)$
has its target in $T(x)$.
There are at least $\delta p$ such arcs.  At most
$\binom{p}{2}$ have both endpoints in $P(x)$, because $G$ is oriented, and at
most $p\bigl(|T(x)|-p\bigr)$ go from $P(x)$ to the rest of the trap.
Consequently, if $p>0$, then
\[
 \delta p
 \le \binom{p}{2}
    +p\bigl(\delta+q(x)-p\bigr).
\]
Dividing by $p$ and rearranging gives $p\le2q(x)-1$.  This is impossible when
$q(x)=0$, so then $p=0$.
\end{proof}

The point of the lemma is that target slack is the only room available for
packing many source vertices that all miss the same target.

\section{Proof of the theorem}

Suppose for contradiction that $G$ has no Seymour vertex.  Then $\delta>0$,
since a vertex of outdegree zero is automatically a Seymour vertex.  Counting
arcs by their sources gives
\[
                         n\delta\le|A(G)|\le\binom n2.
\]
Therefore $n\ge2\delta+1$.  Together with the hypothesis, this leaves
\[
                         c:=n-2\delta\in\{1,2\}.             \tag{2}
\]

Let
\[
                  L=\{v\in V(G):d^+(v)=\delta\},
                  \qquad \ell=|L|.
\]
Some vertex attains the minimum outdegree, so $\ell\ge1$.
Every $v\in L$ is not a Seymour vertex, so $|\NTwo(v)|\le\delta-1$.
Hence
\[
 |U(v)|=n-1-\delta-|\NTwo(v)|\ge n-2\delta=c.               \tag{3}
\]
Thus each minimum-outdegree source contributes at least $c$ unreachable
targets.

Let
\[
 I=\bigl|\{(v,x)\in V(G)^2:x\in U(v)\}\bigr|.
\]
Counting these ordered source--target pairs first by their source and then by
their target gives
\[
                    I=\sum_{v\in V(G)}|U(v)|
                     =\sum_{x\in V(G)}|P(x)|.               \tag{4}
\]
since $x\in U(v)$ and $v\in P(x)$ describe the
same ordered pair $(v,x)$.  From (3),
\[
                              I\ge c\ell>0.                  \tag{5}
\]

Put
\[
                              Q=\sum_{x\in V(G)}q(x).
\]
Since $I>0$, at least one set $P(x)$ is nonempty; the capacity lemma then
implies that at least one $q(x)$ is positive.  Applying the lemma in (4) gives
\[
 I\le\sum_{q(x)>0}\bigl(2q(x)-1\bigr)
   <2\sum_xq(x)=2Q.                                          \tag{6}
\]
So the total target capacity is strictly less than twice the total target
slack.

It remains to compare $Q$ with $\ell$.  Since
$|T(x)|=n-1-d^-(x)$,
\[
 q(x)=n-1-\delta-d^-(x).
\]
Total indegree equals total outdegree, and every vertex outside $L$ has
outdegree at least $\delta+1$.  Therefore
\begin{equation}
\begin{aligned}
 Q
 &=n(n-1-\delta)-\sum_xd^+(x)\\
 &\le n(n-1-\delta)
   -\bigl(\ell\delta+(n-\ell)(\delta+1)\bigr)\\
 &=\ell+n(c-2).
\end{aligned}
\tag{7}
\end{equation}
This degree ledger bounds the total room in all target traps using only the
number of minimum-outdegree vertices.

If $c=1$, then (7) gives
\[
                           0<Q\le\ell-n\le0,
\]
a contradiction.  If $c=2$, then (5)--(7) give
\[
                           I<2Q\le2\ell\le I,
\]
again a contradiction.  Hence $G$ has a Seymour vertex.

\section*{AI contribution and provenance}

Jake Brukhman initiated and directed the investigation, selected its
falsification-first and structural strategy, curated intermediate results,
chose the theorem for publication, and edited the final statement and
exposition.  OpenAI language models---principally a Codex GPT-5-family
multi-agent run and GPT-5.6 Sol---conducted the detailed mathematical
exploration, implemented counterexample searches and verification tools,
discovered the fixed-target capacity argument and its double-counting proof,
and drafted the manuscript.  Claude (Anthropic) was used for an adversarial
audit of an intermediate draft.  Brukhman is responsible for the final
manuscript and its claims.

\begin{thebibliography}{9}

\bibitem{Fisher1996}
D.~C. Fisher,
\emph{Squaring a tournament: a proof of Dean's conjecture},
Journal of Graph Theory \textbf{23} (1996), 43--48.

\bibitem{KanekoLocke2001}
Y.~Kaneko and S.~C. Locke,
\emph{The minimum degree approach for Paul Seymour's distance 2 conjecture},
Congressus Numerantium \textbf{148} (2001), 201--206.

\bibitem{SadhukhanSandeepSen2026}
A.~Sadhukhan, R.~B. Sandeep, and S.~Sen,
\emph{A proof of Seymour's second neighborhood conjecture for oriented graphs
with minimum out-degree equal to 7},
arXiv:2606.30588 (2026).

\bibitem{FidlerYuster2007}
D.~Fidler and R.~Yuster,
\emph{Remarks on the second neighborhood problem},
Journal of Graph Theory \textbf{55} (2007), 208--220.

\bibitem{Ghazal2012}
S.~Ghazal,
\emph{Seymour's second neighborhood conjecture for tournaments missing a
generalized star},
Journal of Graph Theory \textbf{71} (2012), 89--94.

\bibitem{EspunyDiazEtAl2024}
A.~Espuny D\'iaz, A.~Gir\~ao, B.~Granet, and G.~Kronenberg,
\emph{Seymour's second neighbourhood conjecture: random graphs and reductions},
arXiv:2403.02842 (2024).

\bibitem{GuoKangZwaneveld2026}
K.~Guo, R.~J. Kang, and G.~Zwaneveld,
\emph{Seymour-tight orientations},
arXiv:2603.29626 (2026).

\end{thebibliography}

\end{document}
